110.

(a) We may assume $q\in\mathbb{P}$. In Exercise 7.70, let $k=3$ and $x^{(3)}=1^{q}$. Take the scalar product of both sides with $p_n\left(x^{(1)}\right)p_n\left(x^{(2)}\right)$. The right-hand side becomes after some simplification
\[
\mathrm{RHS}=n\sum_{\rho(w)=(n)}q^{\kappa(w(1,2,\ldots,n))}.
\]
The left-hand side becomes (using Corollary 7.21.4)
\[
\mathrm{LHS}=\sum_{\lambda\vdash n}\chi^{\lambda}((n))^2\prod_{u\in\lambda}(q+c(u)).
\]
Now e.g. by the Murnaghan-Nakayama rule we have
\[
\chi^{\lambda}((n))=\begin{cases}(-1)^{n-i}, & \text{if } \lambda=(i,1^{n-i})\\ 0, & \text{otherwise.}\end{cases}
\]
From this we easily get
\[
P_n(q)=\frac{1}{n}\sum_{i=1}^{n}(q+i-1)_n=\frac{1}{n(n+1)}\left((q+n)_{n+1}-(q)_{n+1}\right).
\]
